\chapter{Confidence Intervals}


\begin{definition}[Upper (lower) confidence bound]
	Let $\mathbf X = (X_1, \ldots, X_n)$ be a random sample from a distribution that depends on an unknown parameter $\theta \in \Theta$.
	An upper (lower) $1-\alpha$ confidence bound for $\theta$ is a statistic $U(\mathbf X)$ ($L(\mathbf X)$) such that
	\[
	\mathbb{P}_\theta\left(\theta \leq U(\mathbf X)\right) \geq 1-\alpha
	\qquad
	\left(
	\mathbb{P}_\theta\left(L(\mathbf X) \leq \theta\right) \geq 1-\alpha
	\right),
	\]
	for all $\theta \in \Theta$.
	The corresponding one-sided upper (lower) confidence interval is
	\[
	CI_{\text{upper}} = (-\infty, U(\mathbf X)]
	\qquad
	\left(
	CI_{\text{lower}} = [L(\mathbf X), \infty)
	\right).
	\]
\end{definition}


The quantity $1-\alpha$ is called the confidence level. It describes the probability with which the confidence interval is required to contain the true parameter.
Thus, a confidence bound gives an upper or lower limit for an unknown parameter $\theta$.


If the parameter space $\Theta$ is a strict subset of $\mathbb{R}$, a one-sided confidence interval does not need to extend all the way to $\pm\infty$. Since $\theta \in \Theta$ is already known from the beginning, it is natural to intersect the interval with $\Theta$
\[
CI_{\text{upper}} = \big( -\infty,\, U(\mathbf X) \big] \cap \Theta
\qquad \text{or} \qquad
CI_{\text{lower}} = \big[ L(\mathbf X),\, \infty \big) \cap \Theta,
\]
for the one-sided lower and upper confidence interval, respectively. This intersection does not affect the coverage probability because $\mathbb{P}_\theta\big(\theta \ge L(\mathbf X)\big) \ge 1-\alpha$ and $\mathbb{P}_\theta\big(\theta \le U(\mathbf X)\big) \ge 1-\alpha$ are only required to hold for $\theta \in \Theta$, so excluding values outside $\Theta$ removes no outcome and hence leaves the coverage probability unchanged.

For example, if $\theta$ denotes the rate parameter of an exponential distribution, so that $\Theta = (0, \infty)$, the one-sided upper confidence interval becomes
\[
	CI_{\text{upper}} = \big( 0,\, U(\mathbf X) \big].
\]

\begin{example} \label{example:one sided confidence interval}
	Let $\mathbf X = (X_1, \ldots, X_n)$ be independent and identically exponentially distributed random variables with rate parameter $\lambda \in (0,\infty)$. Thus, the parameter space is
	\(
	\Theta = (0,\infty),
	\)
	since the rate parameter of an exponential distribution must be positive.
	
	To obtain information about the unknown value of $\lambda$, we can construct an upper confidence bound $U(\mathbf X)$.
	
	It is a standard result that
	\begin{align} \label{chi squared distributed}
	2\lambda\sum_{i=1}^n X_i \sim \chi^2_{2n},
	\end{align}
	where $\chi^2_{2n}$ denotes a chi-square distribution with $2n$ degrees of freedom (see for example \cite{bibid} \textcolor{red}{($\star$) }). Therefore,
	\[
	\mathbb{P}\left(
	2\lambda\sum_{i=1}^n X_i
	\leq
	\chi^2_{2n,1-\alpha}
	\right)
	= 1-\alpha,
	\]
	where $\chi^2_{2n,1-\alpha}$ denotes the $(1-\alpha)$-quantile of the chi-square distribution with $2n$ degrees of freedom.
	
	Since $\sum_{i=1}^n X_i > 0$, we can divide both sides of the inequality by $2\sum_{i=1}^n X_i$ and obtain
	\[
	\mathbb{P}\left(
	\lambda
	\leq
	\frac{\chi^2_{2n,1-\alpha}}
	{2\sum_{i=1}^n X_i}
	\right)
	= 1-\alpha.
	\]
	
	Hence, we can define
	\[
	U(\mathbf X)
	=
	\frac{\chi^2_{2n,1-\alpha}}
	{2\sum_{i=1}^n X_i}.
	\]
	This statistic satisfies the definition of an upper $1-\alpha$ confidence bound for $\lambda$. Therefore, the corresponding one-sided upper confidence interval is
	\[
		CI_{\text{upper}}
		=
		(-\infty,U(\mathbf X)]\cap\Theta
		=
		\left(
		-\infty,
		\frac{\chi^2_{2n,1-\alpha}}
		{2\sum_{i=1}^n X_i}
		\right]
		\cap (0,\infty)
		=
		\left(
		0,
		\frac{\chi^2_{2n,1-\alpha}}
		{2\sum_{i=1}^n X_i}
		\right].
	\]
	
	
	Now, we set $\alpha = 0.05, n = 100$ and $\lambda = 2$. We can generate a random realisation of the sum
	\(
	\sum_{i=1}^{100} X_i
	\)
	in R using the code
	
	\begin{verbatim}
		sum(rexp(n = 100, rate = 2))
		[1] 56.55463
	\end{verbatim}
	
	Thus, for this particular realisation, we have
	\(
	\sum_{i=1}^n x_i = 56.55463.
	\)
	
	We can now calculate the upper confidence bound
	\[
	U(\mathbf x)
	=
	\frac{\chi^2_{2\cdot100,\,0.05}}
	{2\cdot56.55463}
	=
	\frac{233.9943}{113.1093}
	=
	2.068745.
	\]
	
	Therefore, the corresponding one-sided $95\%$ confidence interval is
	\[
	CI(\lambda) = (0,2.068745].
	\]
	
	In this particular example, the true parameter value $\lambda=2$ is contained in the confidence interval.
	
	The confidence level of the procedure is $1-\alpha=0.95$. This means that, if we repeatedly generate new samples and construct a confidence interval in the same way, we expect approximately $95\%$ of the resulting intervals to contain the true parameter value $\lambda=2$.
	
	We can illustrate this using a simulation. In Listing~\ref{listing:tbd} \textcolor{red}{($\star$) file: conf\_interval\_exp\_distn.R}, we construct $n_{\text{sim}}=100\,000$ confidence bounds, each based on a different realisation $\mathbf X$ of the sample. For each confidence bound, we check whether the true parameter value $\lambda=2$ is smaller than or equal to the upper confidence bound.
	
	The simulation shows that $94\,908$ out of $100\,000$ confidence bounds are larger than or equal to $\lambda=2$. Thus
	\(
	94.908\%
	\)
	of the resulting one-sided confidence intervals contain the true parameter value. This is close to the confidence level of
	\(
	(1-\alpha)100\% = 95\%
	\).
	Regarding to the law of large numbers, as we increase $n_\text{sim}$, the percentage of confidence intervals which contain the true $\lambda$ should approach $(1 - \alpha)100\%$.
	
\end{example}



If we want to bound the parameter $\theta$ from both sides, we can combine a lower and an upper confidence bound. While a one-sided confidence interval gives a bound for the parameter in one direction, a general confidence interval gives both a lower and an upper limit.

\begin{definition}[Confidence Interval]
	Let $\mathbf X = (X_1, \dots, X_n)$ be a random sample with a distribution depending on an unknown parameter $\theta \in \Theta$. A \emph{confidence interval} for $\theta$ at confidence level $1-\alpha$, where $\alpha \in (0,1)$, is an interval
	\[
	CI = [L(\mathbf X), U(\mathbf X)],
	\]
	where $L(\mathbf X)$ and $U(\mathbf X)$ are statistics satisfying
	\[
	L(\mathbf X) \leq U(\mathbf X)
	\]
	and
	\[
	\mathbb{P}_\theta\big(
	L(\mathbf X) \leq \theta \leq U(\mathbf X)
	\big)
	\geq 1-\alpha
	\]
	for all $\theta \in \Theta$.
\end{definition}


\begin{example}
For the scenario from Example~\ref{example:one sided confidence interval}, we can also construct a two-sided confidence interval
\[
	CI(\lambda) = [L(\mathbf X), U(\mathbf X)].
\]
Again, we use the fact that $2\lambda\sum_{i=1}^n X_i$ is chi squared distributed as we saw in \eqref{chi squared distributed}
Therefore,
\begin{align*}
	\mathbb{P}\left(
	\chi^2_{2n,\frac{\alpha}{2}}
	\leq
	2\lambda\sum_{i=1}^n X_i
	\leq
	\chi^2_{2n,1-\frac{\alpha}{2}}
	\right)
	&=
	\mathbb{P}\left(
	2\lambda\sum_{i=1}^n X_i
	\leq
	\chi^2_{2n,1-\frac{\alpha}{2}}
	\right)
	-
	\mathbb{P}\left(
	2\lambda\sum_{i=1}^n X_i
	<
	\chi^2_{2n,\frac{\alpha}{2}}
	\right) 
	\\
	&=
	\left(1-\frac{\alpha}{2}\right) - \frac{\alpha}{2}
	\\
	&=
	1-\alpha.
\end{align*}

Since
\(
2\sum_{i=1}^n X_i > 0,
\)
we can divide all parts of the inequality by it and this gives
\[
	\mathbb{P}\left(
	\frac{\chi^2_{2n,\frac{\alpha}{2}}}
	{2\sum_{i=1}^n X_i}
	\leq
	\lambda
	\leq
	\frac{\chi^2_{2n,1-\frac{\alpha}{2}}}
	{2\sum_{i=1}^n X_i}
	\right)
	=
	1-\alpha.
\]

Thus, the lower and upper confidence bounds can be defined as
\[
	L(\mathbf X)
	=
	\frac{\chi^2_{2n,\frac{\alpha}{2}}}
	{2\sum_{i=1}^n X_i}
	\qquad\text{and}\qquad
	U(\mathbf X)
	=
	\frac{\chi^2_{2n,1-\frac{\alpha}{2}}}
	{2\sum_{i=1}^n X_i}.
\]

Consequently, a $(1-\alpha)$ confidence interval for $\lambda$ is given by
\[
	CI(\lambda)
	=
	\left[
	\frac{\chi^2_{2n,\frac{\alpha}{2}}}
	{2\sum_{i=1}^n X_i},
	\frac{\chi^2_{2n,1-\frac{\alpha}{2}}}
	{2\sum_{i=1}^n X_i}
	\right].
\]

For the data considered in Example~\ref{example:one sided confidence interval}, we have $\alpha=0.05$, $n=100$, $\lambda=2$
and the observed value is
\(
	\sum_{i=1}^{100}x_i=56.55463
\).
The resulting $95\%$ confidence interval is given by
\[
	CI(\lambda)
	=
	\left[
	\frac{\chi^2_{2\cdot100,0.025}}
	{2\cdot56.55463},
	\frac{\chi^2_{2\cdot100,0.975}}
	{2\cdot56.55463}
	\right]
	=
	[1.44, 2.13].
\]


In contrast to the one-sided confidence interval, which only provides an upper bound for $\lambda$, the two-sided confidence interval provides both a lower and an upper bound. This is useful when we do not have a specific reason to expect that $\lambda$ is only of concern in one direction. The two-sided confidence interval therefore gives information about whether $\lambda$ could be  smaller or larger than the observed estimate.

Another difference is how the error probability $\alpha$ is distributed.
For the one-sided confidence interval, the whole error probability $\alpha$
is placed in one tail of the distribution. In
Example~\ref{example:one sided confidence interval}, it is placed on the
right side. The interval fails exactly when the true parameter $\lambda$
lies to the right of the upper limit. For the two-sided confidence interval,
the error probability is split between both tails, with $\frac{\alpha}{2}$ in each tail.

We can check this with the simulation of Listing~\ref{todo}
\textcolor{red}{($\star$) file conf\_interval\_exp\_distn.R}. For each
simulated sample, we determine whether the true parameter is to the left of the
interval, inside the interval, or to the right of the interval. The results are
as follows. In $94.956\%$ of the simulations, the confidence interval contains
the true parameter $\lambda = 2$. This is close to the confidence level
$(1 - \alpha) 100\% = 95\%$. The true parameter lies to the left of the interval in
$2.526\%$ of the simulations and to the right of the interval in $2.518\%$.
Both values are close to $\frac{\alpha}{2} = 2.5\%$. So the error probability is
split almost equally between the two sides.

This aligns with our expectation because the probability that the true parameter $\lambda$ is left to the interval is given by
\[
\mathbb{P} (\lambda < L \left( \mathbf X \right))
= \mathbb{P} \left( 2\lambda \sum_{i = 1}^n X_i < \chi^2_{2n;\,\frac{\alpha}{2}} \right)
= \frac{\alpha}{2},
\]
and the probability that $\lambda$ is to the right of the interval is
\[
\mathbb{P}(\lambda > U(\mathbf X))
= \mathbb{P} \left( 2\lambda \sum_{i = 1}^n X_i > \chi^2_{2n;\,1- \frac{\alpha}{2}} \right)
= 1 - \left( 1 - \frac{\alpha}{2} \right)
= \frac{\alpha}{2}.
\]

Instead of the split $\alpha = \frac{\alpha}{2} + \frac{\alpha}{2}$, we can use any
split $\alpha = \alpha_1 + \alpha_2$ with $\alpha_1, \alpha_2 > 0$. Then
\[
	\left[
		\frac{\chi^2_{2n;\,\alpha_1}}{2 \sum_{i = 1}^n X_i} ,\;
		\frac{\chi^2_{2n;\,1 - \alpha_2}}{2 \sum_{i = 1}^n X_i}
	\right]
\]
is also a confidence interval for $\lambda$ with confidence level $1 - \alpha$.

\end{example}



\section{Exact vs. Asymptotic CIs}

\textcolor{red}{($\star$) Exact vs Asymptotic CI, plot x = n, y = P(noncoverage), dann müsste man bei exact sehen, dass n keinen einfluss auf p(nc) hat und bei asymptotischen CI sieht man dass diese mit größeren n besser wird, gib kleines Beispiel und plots (siehe \texttt{exact\_vs\_asymptotic.R})}

\begin{definition}[Exact confidence interval]
	A confidence interval $CI$ for a parameter $\theta$ is called \textbf{exact}, if and only if
	\[
		\mathbb{P}_{\theta} \left( \theta \in CI \right) = 1 - \alpha.
	\]
\end{definition}

\begin{definition}[Asymptotic confidence interval]
	A confidence interval $CI$ for a parameter $\theta$ is called \textbf{asymptotic}, if
	\[
	\mathbb{P}_{\theta} \left( \theta \in CI \right) \overset{n \to \infty}{\longrightarrow} 1 - \alpha.
	\]
	where $n$ is the sample size.
\end{definition}

So the main difference between exact and asymtotic confidence intervals is, that exact confidence intervals guarantee a confidence niveau of $1 - \alpha$ even for a small sample size $n$, but asymptotic confidence intervals only have confidence niveau $1 - \alpha$ in the limit $n \to \infty$.

	

\begin{example}[Exact vs. Asymptotic Confidence Interval]
	Let $X_1, \ldots , X_n \overset{\text{iid}}{\sim} \mathcal N (\mu, \sigma)$. In the following, we want to construct a confidence interval for the mean $\mu$. Since the mean of the normal distribution can be any real number, the parameter space is $\theta = \mathbb R$. Consider two cases.
	
	
	\textbf{Case 1: Standard deviation $\sigma$ is known} \\
	In this case, we can construct an exact confidence interval the following way.
	Observe that the arithmetic mean of the observations fulfils $\bar X = \frac{1}{n} \sum_{i=1}^n X_i \sim \mathcal N \left( \mu, \frac{\sigma^2}{n} \right)$ and by scaling this term, we get
	\[
		\frac{\sqrt{n}(\bar X - \mu)}{\sigma} \sim \mathcal N (0, 1).
	\]
	This distribution holds for every sample size $n$ so it is not only asymptotic. 
	Let $\alpha \in (0, 1)$ be arbitrary. We can construct a confidence interval by 
	\begin{align}
		& 1 - \alpha \notag \\
		=&
		\left( 1 - \frac{\alpha}{2} \right) - \frac{\alpha}{2} \notag \\
		=&
		\mathbb{P} \left(  
		\frac{\sqrt{n}(\bar X - \mu)}{\sigma}
		\leq z_{1 - \frac{\alpha}{2}} \right)
		-
		\mathbb{P} \left( 
		\frac{\sqrt{n}(\bar X - \mu)}{\sigma}
		< z_{\frac{\alpha}{2}} \right) \notag \\
		=&
		\mathbb{P} \left( z_{\frac{\alpha}{2}} \leq 
		\frac{\sqrt{n}(\bar X - \mu)}{\sigma}
		\leq z_{1 - \frac{\alpha}{2}} \right) \notag \\
		=&
		\mathbb{P} 
		\left( 
		\bar X - \frac{\sigma}{\sqrt{n}} z_{\frac{\alpha}{2}} \geq  \mu \geq \bar X - \frac{\sigma}{\sqrt{n}} z_{1 - \frac{\alpha}{2}}		
		\right)
	\end{align}
	For the quantiles of the normal distributino it holds $z_p = - z_{1-p}$ for all $p \in (0, 1)$, so we can substitute the quantile $z_{\frac{\alpha}{2}}$ above. The $1- \alpha$ confidence interval is given by
	\[
		CI_{1 - \alpha}^{\sigma \text{ known}} (\mu) 
		=
		\left[ 
		\bar X - \frac{\sigma}{\sqrt{n}} z_{1 - \frac{\alpha}{2}}
		,
		\bar X + \frac{\sigma}{\sqrt{n}} z_{1 - \frac{\alpha}{2}}
		\right]
	\]
	We have shown that the probability that the real parameter $\mu$ is in $CI_{1 - \alpha}^{\sigma \text{ known}} (\mu)$ is $1 - \alpha$ independently of the sample size $n$, so $CI_{1 - \alpha}^{\sigma \text{ known}} (\mu)$ is an exact confidence interval. 
	
	
	\textbf{Case 2: Standard deviation $\sigma$ is unknown} \\
	
	
	From the Central Limit Theorem, we know that
	\[
		\frac{\sqrt{n} (\bar X - \mu)}{\sigma} \overset{n \rightarrow \infty}{\longrightarrow} \mathcal N (0, 1).
	\]
	Since $\sigma$ is unknown, we need to substitute it somehow.
	It holds 
	\[
		S_n := \sqrt{\frac{1}{n-1} \sum_{i=1}^n \left( X_i - \bar X \right)^2 }  \overset{p}{\longrightarrow} \sigma 
	\]
	\textcolor{red}{($\star$) Warum ist $S_n$ so verteilt?}
	
	With Slutsky's Theorem, we can conclude that 
	\[
		\frac{\sqrt{n} (\bar X - \mu)}{S_n} \overset{n \rightarrow \infty}{\longrightarrow} \mathcal N (0, 1).
	\]
	and with that knowledge we can create a $1 - \alpha$ confidence interval similarily to case 1
	\[
		CI_{1 - \alpha}^{\sigma \text{ unknown}} (\mu) 
		=
		\left[ 
		\bar X - \frac{S_n}{\sqrt{n}} z_{1 - \frac{\alpha}{2}}
		,
		\bar X + \frac{S_n}{\sqrt{n}} z_{1 - \frac{\alpha}{2}}
		\right].
	\]
	This confidence interval is only asymptotic because \textcolor{red}{($\star$) Warum nur asymptotic? Warum funktioniert die einsetzung mit slutsky nur asymptotic?}
	
	
	In Figure~\ref{fig:exact_vs_asymptotic}, we can see the non-coverage probability of the intervals $CI_{1 - \alpha}^{\sigma \text{ known}} (\mu) $ and $CI_{1 - \alpha}^{\sigma \text{ unknown}} (\mu) $ dependent on the sample size. In Listing~\ref{listing:exact vs asymptotic} is the code used for generating Figure~\ref{fig:exact_vs_asymptotic}. 
	
	We can observe that $CI_{1 - \alpha}^{\sigma \text{ known}} (\mu) $ maintains a non-coverage probability close to \(\alpha=0.1\) for all sample sizes $n$, while the non-coverage probability of $CI_{1 - \alpha}^{\sigma \text{ unknown}} (\mu) $ approaches \(\alpha\) as \(n\) increases.
	
	\begin{figure}[htbp]
		\centering
		\includegraphics[width=0.6\textwidth]{figures/exact\_vs\_asymptotic\_CI.png}
		\caption{Non-coverage probability of $CI_{1 - \alpha}^{\sigma \text{ known}} (\mu)$ and $CI_{1 - \alpha}^{\sigma \text{ unknown}} (\mu)$ based on 1000 Monte Carlo simulations for each sample size $n=2,\ldots,100$. For each simulation, a sample from \(N(\mu,\sigma^2)\) with $\mu=10$ and $\sigma=100$ is generated, and it is recorded whether \(\mu\) is contained in the respective confidence interval.}
		\label{fig:exact_vs_asymptotic}
	\end{figure}
		
\end{example}


